\chapter{Normalization and weak adjoint measures}
\label{chap:normalization}

The scalar multiplying a prescribed Monge--Amp\`ere density is unique even when the potential is not. Varying the logarithmic density therefore gives a well-defined functional. Mixed products and domination make it concave; its supporting functionals determine probability measures. Continuity of the scalar then controls their concentration and produces a dual variational problem. In the minimal singularity type, these probabilities also identify the potential selected by a vanishing exponential parameter.

Fix a connected compact complex manifold $X$ of dimension $n$ and a smooth positive volume form $\mathrm{d}V$. We first take $\theta$ smooth, semipositive and big, and a model reference $\varphi$ satisfying \eqref{eq:als-model-condition}. The equation is supplied by \cref{thm:als-prescribed-lp}. Function spaces are real, and their norms use $\mathrm{d}V$. The final selection argument will allow an arbitrary smooth big background.

\section{The scalar functional}\label{sec:normalization-scalar}

Changing the density changes the scalar in the equation. Comparison determines the direction of this change, and multiplication of the density by a constant gives an exact translation rule.

For a bounded real measurable function $F$, \cref{thm:als-prescribed-lp} gives $u\sim\varphi$ and a unique $c_\varphi(F)>0$ such that
\begin{equation}\label{eq:relative-normalization}
 \mathrm{MA}^{\mathrm{ext}}_\theta(u)
       =c_\varphi(F)\mathrm{e}^F\,\mathrm{d}V.
\end{equation}
Define
\[
 B_\varphi(F)=\log c_\varphi(F).
\]
Functions equal $\mathrm{d}V$-almost everywhere are identified. Scalar uniqueness makes $B_\varphi$ well defined and attaches it to the prescribed ordinary type, even when the normalized potential is not unique.

\Needspace{13\baselineskip}
\begin{proposition}\label{prop:relative-normalization-order}
Let $F,G\in L^\infty(\mathrm{d}V)$ be real. The functional $B_\varphi$ is decreasing, and for every constant $a\in\R$,
\[
 B_\varphi(F+a)=B_\varphi(F)-a.
\]
Consequently
\begin{equation}\label{eq:relative-normalization-lipschitz}
 -\mathop{\rm ess\,sup}_X(G-F)
 \leq B_\varphi(G)-B_\varphi(F)
 \leq-\mathop{\rm ess\,inf}_X(G-F),
\end{equation}
and $|B_\varphi(G)-B_\varphi(F)|\leq\|G-F\|_\infty$.
\end{proposition}
\begin{proof}
Suppose $F\leq G$ almost everywhere and, to obtain a contradiction, $c_\varphi(F)<c_\varphi(G)$. Let $u,v\sim\varphi$ solve the two equations. Then, for every constant $a$,
\[
 \mathrm{MA}^{\mathrm{ext}}_\theta(u)
 \leq\frac{c_\varphi(F)}{c_\varphi(G)}
             \mathrm{MA}^{\mathrm{ext}}_\theta(v+a).
\]
The ratio is strictly less than one. Bounded equivalence transfers the genuine smaller-background envelope from $\varphi$ to $u$, and $v+a\preceq u$. Singular strict domination gives $u\geq v+a$ for every $a$. This is impossible at any point where both potentials are finite. Thus $B_\varphi(F)\geq B_\varphi(G)$.

Replacing $F$ by $F+a$ and the constant by $\mathrm{e}^{-a}c_\varphi(F)$ leaves the right side of \eqref{eq:relative-normalization} unchanged. Scalar uniqueness proves the translation rule. Apply monotonicity to the two constant translates of $F$ bounding $G$ to obtain \eqref{eq:relative-normalization-lipschitz}.
\end{proof}

For a smooth real $h$, replacing $\theta$ by $\theta+\ddc h$ and each potential by $u-h$ preserves the current and the prescribed ordinary type. Transporting the solution problem in this way leaves its scalar unchanged.

\section{Concavity from mixed products}

A convex combination of two solutions stays in the prescribed ordinary type. The mixed density inequality of \cref{lem:singular-convex-density} bounds its measure from below; strict domination then compares that bound with the scalar for the interpolated density.

\begin{theorem}\label{thm:relative-normalization-concavity}
The functional $B_\varphi$ is concave on the real space $L^\infty(\mathrm{d}V)$: for $F,G$ in that space,
\[
 B_\varphi(tF+(1-t)G)
 \geq tB_\varphi(F)+(1-t)B_\varphi(G),\qquad0\leq t\leq1.
\]
\end{theorem}
\begin{proof}
Let $u,v\sim\varphi$ solve the equations for $F,G$. Put
$H=tF+(1-t)G$ and $D=c_\varphi(F)^t c_\varphi(G)^{1-t}$. \Cref{lem:singular-convex-density} and the arithmetic--geometric mean inequality give
\[
 \begin{split}
 \mathrm{MA}^{\mathrm{ext}}_\theta(tu+(1-t)v)
 &\geq\bigl(t(c_\varphi(F)\mathrm{e}^F)^{1/n}
       +(1-t)(c_\varphi(G)\mathrm{e}^G)^{1/n}\bigr)^n\,\mathrm{d}V\\
 &\geq D \mathrm{e}^H\,\mathrm{d}V.
 \end{split}
\]
The averaged potential has ordinary type $\varphi$. If $c_\varphi(H)<D$, its measure dominates that of a solution for $H$ by a factor strictly greater than one. Apply singular strict domination to that solution and every constant translate of the averaged potential, as in the proof of \cref{prop:relative-normalization-order}. The same contradiction gives $c_\varphi(H)\geq D$. Taking logarithms proves concavity.
\end{proof}

For smooth Hermitian data, this concavity agrees with the differential calculation of \cref{thm:b-derivatives}.

\section{Continuity for bounded densities}

To rule out concentration of a supporting measure on a set of zero volume, we need continuity under bounded changes of density that tend to zero in $L^1$. The relative estimate gives common barriers for the solutions, and their common measure bound identifies every limiting scalar.

\begin{theorem}\label{thm:normalization-bounded-lone}
If $F_j,F$ are uniformly bounded and $F_j\to F$ in $L^1(\mathrm{d}V)$, then
\[
 B_\varphi(F_j)\longrightarrow B_\varphi(F).
\]
\end{theorem}
\begin{proof}
Choose solutions $u_j\sim\varphi$ normalized by $\sup_Xu_j=0$, and write $c_j=c_\varphi(F_j)$. If $|F_j|\leq M$, order and translation give
\[
 \mathrm{e}^{-M}c_\varphi(0)\leq c_j\leq \mathrm{e}^M c_\varphi(0).
\]
In particular their Monge--Amp\`ere measures are dominated by a fixed multiple of $\mathrm{d}V$. \Cref{thm:relative-linfty} gives a constant $C$ such that
\[
 \varphi-C\leq u_j\leq\varphi.
\]
From any subsequence extract $c_j\to c>0$ and $u_j\to u$ in $L^1(X)$. \Cref{cor:common-barrier-continuity} gives convergence in capacity, $u\sim\varphi$, and raw weak convergence of their measures. On the other hand, boundedness makes the exponential Lipschitz on the common range, so $c_j\mathrm{e}^{F_j}\to c\mathrm{e}^F$ in $L^1(\mathrm{d}V)$. Therefore
\[
 \mathrm{MA}^{\mathrm{ext}}_\theta(u)=c\mathrm{e}^F\,\mathrm{d}V.
\]
Scalar uniqueness gives $c=c_\varphi(F)$. Every subsequence has a further subsequence with this scalar limit, proving convergence of the full sequence and then of its logarithms.
\end{proof}

\section{Supporting probability measures}\label{sec:normalization-supports}

A supporting functional records the first-order response of the scalar to a change of logarithmic density. Monotonicity and the translation rule force its negative to be integration against a probability. Work first on $C(X)$ with the uniform norm. For $F\in C(X)$, let $\mathcal A_\varphi(F)$ be the set of Borel probability measures $\nu$ satisfying
\begin{equation}\label{eq:adjoint-probability-set}
 B_\varphi(G)\leq B_\varphi(F)-\int_X(G-F)\,\mathrm{d}\nu
 \qquad\text{for every }G\in C(X).
\end{equation}
Thus $-\nu$ is a supergradient of $B_\varphi$ at $F$.

\begin{theorem}\label{thm:weak-adjoint-probabilities}
For every $F\in C(X)$, the set $\mathcal A_\varphi(F)$ is nonempty, convex and weakly compact. Every one of its measures is absolutely continuous with respect to $\mathrm{d}V$.
\end{theorem}
\begin{proof}
Write $C_\varphi=-B_\varphi$, a convex Lipschitz function. For fixed $F\in C(X,\mathbb R)$ and a direction $h\in C(X,\mathbb R)$, the limit
\[
 p(h)=\lim_{t\downarrow0}
        \frac{C_\varphi(F+th)-C_\varphi(F)}{t}
\]
exists and is finite: convexity makes the difference quotients monotone in $t>0$, and their absolute values are at most $\|h\|_\infty$. Positive homogeneity is immediate. Applying convexity with another direction $k\in C(X,\mathbb R)$ at the midpoint of $F+2th$ and $F+2tk$ and then letting $t\downarrow0$ gives $p(h+k)\leq p(h)+p(k)$. Thus $p$ is sublinear.

Hahn--Banach extends the zero functional on $\{0\}$ to a linear functional $\ell\leq p$. Since $p(h)\leq\|h\|_\infty$, application to $h$ and $-h$ makes $\ell$ continuous. The monotonicity of the convex difference quotients gives
\[
 C_\varphi(F+h)\geq C_\varphi(F)+\ell(h).
\]
The translation identity, tested with positive and negative constants, implies $\ell(1)=1$. If $h\geq0$, monotonicity gives
$C_\varphi(F-th)\leq C_\varphi(F)$ for $t>0$, whereas the supporting inequality gives the reverse lower bound $C_\varphi(F)-t\ell(h)$. Hence $\ell(h)\geq0$. By the Riesz representation theorem, $\ell(h)=\int_Xh\,\mathrm{d}\nu$ for a positive probability measure $\nu$. This is exactly \eqref{eq:adjoint-probability-set}, proving nonemptiness. Each inequality in that definition specifies a weakly closed half-space in the compact space of probability measures. Their intersection is closed and convex, hence weakly compact.

It remains to prove absolute continuity. Let $K$ be compact with $\mathrm{d}V(K)=0$. Choose continuous functions $0\leq\chi_j\leq1$ equal to one on $K$, whose supports have volumes tending to zero. For any $\nu\in\mathcal A_\varphi(F)$, the supporting inequality with $G=F+\chi_j$ gives
\[
 0\leq\nu(K)\leq\int_X\chi_j\,\mathrm{d}\nu
 \leq B_\varphi(F)-B_\varphi(F+\chi_j).
\]
The right side tends to zero by \cref{thm:normalization-bounded-lone}. Thus $\nu(K)=0$. Inner regularity of a finite Borel measure on the compact manifold proves the same for every Borel set of zero volume. The Radon--Nikodym theorem now gives $\mathrm{d}\nu=h_\nu \mathrm{d}V$ with $h_\nu\geq0$ and $\int_Xh_\nu \mathrm{d}V=1$.
\end{proof}

\begin{proposition}
\label{prop:weak-adjoint-stability}
If $F_j\to F$ uniformly and $\nu_j\in\mathcal A_\varphi(F_j)$, every weak cluster point of $\nu_j$ belongs to $\mathcal A_\varphi(F)$. In particular, if the latter set consists of a single measure, $\nu_j$ converges to it.
\end{proposition}
\begin{proof}
Take a weakly convergent subsequence and fix $G\in C(X)$. In
\[
 B_\varphi(G)\leq B_\varphi(F_j)-\int_X(G-F_j)\,\mathrm{d}\nu_j,
\]
uniform continuity of $B_\varphi$, uniform convergence of $F_j$ and weak convergence of the probabilities allow passage to the limit. The resulting inequality holds for every $G$, so the limiting measure is in $\mathcal A_\varphi(F)$. If that set is a singleton, every convergent subsequence has the same limit; compactness proves the final assertion.
\end{proof}

\section{Directional derivatives and the smooth identification}

The range of integrals against the supporting probabilities is exactly the gap between the two directional derivatives. A singleton therefore identifies a derivative, and the smooth equation identifies that measure with the Gauduchon adjoint. For $F,G\in C(X)$, write $D^+B_\varphi(F)[G]$ and $D^-B_\varphi(F)[G]$ for the limits of the difference quotient as $t\downarrow0$ and $t\uparrow0$, respectively.

\begin{proposition}\label{prop:normalization-directional}
The one-sided derivatives exist, and
\begin{align*}
 D^+B_\varphi(F)[G]&=-\max_{\nu\in\mathcal A_\varphi(F)}\int_XG\,\mathrm{d}\nu,\\
 D^-B_\varphi(F)[G]&=-\min_{\nu\in\mathcal A_\varphi(F)}\int_XG\,\mathrm{d}\nu.
\end{align*}
Consequently $B_\varphi$ is G\^ateaux differentiable at $F$ if and only if $\mathcal A_\varphi(F)$ is a singleton. Its derivative is then $G\mapsto-\int_XG\,\mathrm{d}\nu$.
\end{proposition}
\begin{proof}
Existence and finiteness of the one-sided limits follow from concavity and the Lipschitz bound. For $t>0$, every supporting measure at $F$ gives
\[
 \frac{B_\varphi(F+tG)-B_\varphi(F)}{t}
 \leq-\int_XG\,\mathrm{d}\nu.
\]
Conversely choose $\nu_t\in\mathcal A_\varphi(F+tG)$ and apply its supporting inequality to $F$. This gives the lower bound $-\int_XG\,\mathrm{d}\nu_t$ for the same quotient. Along a sequence $t\downarrow0$, extract a weak limit of $\nu_t$. It belongs to $\mathcal A_\varphi(F)$ by \cref{prop:weak-adjoint-stability}. The two bounds identify the right derivative with the negative maximum. Apply this identity to $-G$ to obtain the left derivative.

If the set is a singleton, the two derivatives coincide for every $G$ and equal the displayed continuous linear functional. Conversely, equality of the two one-sided derivatives for every $G$ forces all supporting probabilities to have identical integrals against every continuous function. The Riesz theorem makes them identical.
\end{proof}

\begin{theorem}
\label{thm:normalization-gauduchon-recovery}
Let $\theta=\omega$ be Hermitian, let $\varphi=0$, and let $F$ be smooth. If $u$ is the smooth normalized solution of \eqref{eq:relative-normalization}, put $\omega_u=\omega+\ddc u$. Let $\rho_u$ be the positive smooth Gauduchon factor with $\ddc(\rho_u\omega_u^{n-1})=0$ and $\int_X\rho_u\omega_u^n=1$, and put $\mathrm{d}\mu_u=\rho_u\omega_u^n$. Then
\[
 \mathcal A_0(F)=\{\mu_u\}.
\]
For $n=1$ the same assertion uses $\mu_u=\omega_u/\int_X\omega_u$.
\end{theorem}
\begin{proof}
In dimensions at least two, \cref{thm:b-derivatives} shows that for every smooth real $G$,
\[
 \left.\frac{\mathrm{d}}{\mathrm{d}t}\right|_{t=0}B_0(F+tG)
 =-\int_XG\,\mathrm{d}\mu_u.
\]
Scalar uniqueness identifies its normalizing constant with the present $B_0$. For any $\nu\in\mathcal A_0(F)$, testing the supporting inequality at $F+tG$ for both signs of $t$ gives equality of the integrals of $G$ against $\nu$ and $\mu_u$. Smooth functions are uniformly dense in $C(X)$, so $\nu=\mu_u$. Nonemptiness proves the assertion. In dimension one Stokes gives
$c_0(F)=\int_X\omega/\int_X\mathrm{e}^F \mathrm{d}V$, whose logarithmic derivative gives the stated measure directly.
\end{proof}

Thus for a continuous $F$ in the Hermitian case, smooth approximations $F_j\to F$ uniformly produce Gauduchon adjoint probabilities whose weak cluster points all belong to $\mathcal A_0(F)$ and are absolutely continuous. If $\mathcal A_0(F)$ is a singleton, the weak limit is independent of this approximation. This is a conditional answer to the question in \cref{chap:gauduchon}, for approximations controlled through the logarithmic density.

In dimension $n\geq2$, the Hermitian zero-reference case has a quantitative strengthening: \cref{cor:gaud-adjoint-lq} gives uniform $L^q$ bounds for all these densities when $\|F\|_\infty$ is bounded and $1<q<n/(n-1)$. It includes $q=n/(n-1)$ under the positive lower-volume hypothesis. Its proof uses the finite Gauduchon-current formulation of \cref{chap:gauduchon-variational-volume}; no such bound for an arbitrary singular reference is inferred here.

Let $\mathcal D_\omega^0$ be the closure of bounded density roots in the norm of \cref{subsec:bounded-root-approximation}. For an integrable nonnegative density $f$ whose root belongs to $\mathcal D_\omega^0$, \cref{prop:integrable-threshold-adjoints} will give a similar representation for $F\mapsto\log b_\omega(\mathrm{e}^Ff)$. That functional is defined by a subsolution threshold; its construction does not require an exact equation for $f$ or prescribe a singularity type.

\section{A convex functional on probability measures}

The scalar functional also has a dual description. This gives a variational problem in the prescribed density and a useful compactness statement for its supporting measures. The variable in this problem is a probability measure.

Let $\mathcal P(X)$ denote the Borel probability measures on $X$. Define
\begin{equation}\label{eq:normalization-dual-definition}
 H_\varphi(\nu)
 =\sup_{G\in C(X)}\left(B_\varphi(G)+\int_XG\,\mathrm{d}\nu\right),
 \qquad \nu\in\mathcal P(X).
\end{equation}
The value $+\infty$ is allowed. Taking $G=0$ shows that
$H_\varphi\geq B_\varphi(0)$, so the value $-\infty$ never occurs.

\begin{proposition}\label{prop:normalization-dual}
The functional $H_\varphi$ is convex and weakly lower semicontinuous. For every $F\in C(X)$,
\begin{equation}\label{eq:normalization-dual-minimum}
 B_\varphi(F)
 =\min_{\nu\in\mathcal P(X)}
       \left(H_\varphi(\nu)-\int_XF\,\mathrm{d}\nu\right).
\end{equation}
The minimizers are exactly $\mathcal A_\varphi(F)$.
\end{proposition}
\begin{proof}
Every function of $\nu$ in the supremum defining $H_\varphi$ is affine and weakly continuous. Their supremum is convex and lower semicontinuous. The definition gives
$H_\varphi(\nu)-\int_XF\,\mathrm{d}\nu\geq B_\varphi(F)$.
Equality holds precisely when
\[
 B_\varphi(G)+\int_XG\,\mathrm{d}\nu
 \leq B_\varphi(F)+\int_XF\,\mathrm{d}\nu
 \quad\text{for every }G\in C(X).
\]
This is the defining inequality for $\mathcal A_\varphi(F)$. Its nonemptiness, proved in \cref{thm:weak-adjoint-probabilities}, proves attainment.
\end{proof}

\Needspace{14\baselineskip}
\begin{theorem}\label{thm:normalization-dual-compactness}
For every finite real $L$, the set
\[
 \mathcal K_L=\{\nu\in\mathcal P(X):H_\varphi(\nu)\leq L\}
\]
is weakly compact. Its measures are uniformly absolutely continuous with respect to $\mathrm{d}V$: for every $\epsilon>0$ there is a $\delta>0$ such that
\begin{equation}\label{eq:normalization-dual-uniform-ac}
 \mathrm{d}V(E)<\delta\quad\Longrightarrow\quad
 \sup_{\nu\in\mathcal K_L}\nu(E)<\epsilon
\end{equation}
for every Borel set $E$. In particular, writing $\mathrm{d}\nu=h_\nu \mathrm{d}V$, the densities in this set are uniformly integrable in $L^1(\mathrm{d}V)$. On $\mathcal K_L$, weak convergence of measures also implies weak convergence of their densities in $L^1(\mathrm{d}V)$.
\end{theorem}
\begin{proof}
Weak compactness follows from lower semicontinuity and compactness of $\mathcal P(X)$. The empty sublevel is harmless; otherwise $L\geq B_\varphi(0)$.
Fix $t>0$. By \cref{thm:normalization-bounded-lone} there is $\delta_t>0$ such that
\[
 0\leq\chi\leq1,\quad \chi\in C(X),\quad
 \int_X\chi\,\mathrm{d}V<\delta_t
 \quad\Longrightarrow\quad
 B_\varphi(t\chi)>B_\varphi(0)-1.
\]
Indeed, failure would give a sequence of such $\chi$ tending to zero in $L^1$, contrary to that theorem. For $\nu\in\mathcal K_L$, the defining supremum gives
\[
 t\int_X\chi\,\mathrm{d}\nu
 \leq L-B_\varphi(t\chi)
 <L-B_\varphi(0)+1.
\]
Choose $t$ so large that the last numerator divided by $t$ is less than $\epsilon$. If $\mathrm{d}V(E)<\delta_t/2$, outer regularity supplies an open neighborhood $O$ of $E$ with $\mathrm{d}V(O)<\delta_t$. For any compact $K\subset E$, choose a continuous cutoff equal to one on $K$, between zero and one, and supported in $O$. The last inequality gives $\nu(K)<\epsilon$, uniformly in $\nu$. Inner regularity gives the same bound for $E$, with a non-strict inequality. Choosing $t$ initially with bound $\epsilon/2$ proves \eqref{eq:normalization-dual-uniform-ac} as stated.

This proves absolute continuity of every measure in the sublevel, and the uniform small-set criterion for the densities, whose integrals all equal one. These two facts are exactly uniform integrability on a finite measure space.

Finally suppose $\nu_j\to\nu$ weakly within the sublevel. For a bounded measurable $G$, choose continuous $G_k$ with
$|G_k|\leq\|G\|_\infty$ and $G_k\to G$ in $\mathrm{d}V$-measure. Such approximations follow, for example, by approximating simple functions using regularity and continuous cutoffs. Uniform absolute continuity gives
\[
 \sup_{\sigma\in\mathcal K_L}\int_X|G_k-G|\,\mathrm{d}\sigma\longrightarrow0.
\]
To see this, split into $\{|G_k-G|\leq a\}$ and its complement; the first integral is at most $a$, and the second is bounded by $2\|G\|_\infty$ times a uniformly small measure. For fixed $k$, weak convergence applies to $G_k$. Passing first $j\to\infty$ and then $k\to\infty$ proves convergence against $G$. This is precisely weak $L^1(\mathrm{d}V)$ convergence of the densities. In particular every sequence in the sublevel has a weakly convergent $L^1$ subsequence.
\end{proof}

This compactness is obtained from the continuity of the equation's normalizing constant. It uses no uniform bound on Monge--Amp\`ere masses of arbitrary bounded potentials. It does not give a common $L^p$ bound for the densities, for any $p>1$.

\section{Bounded measurable densities and stability}

The preceding construction extends to the bounded measurable functions on which $B_\varphi$ was originally defined. No finitely additive measures are needed.

\begin{theorem}\label{thm:normalization-measurable-adjoint}
For every $F\in L^\infty(\mathrm{d}V)$ there is a nonempty set of absolutely continuous probability measures, again denoted $\mathcal A_\varphi(F)$, characterized by
\begin{equation}\label{eq:measurable-adjoint-support}
 B_\varphi(G)\leq B_\varphi(F)-\int_X(G-F)\,\mathrm{d}\nu
 \qquad\text{for every }G\in L^\infty(\mathrm{d}V).
\end{equation}
For continuous $F$, this agrees with the preceding definition. Formula \eqref{eq:normalization-dual-minimum} holds for bounded measurable $F$ as well, with these measures as its minimizers. In that formula the objective is declared to be $+\infty$ whenever $H_\varphi(\nu)=+\infty$; otherwise $\nu$ is absolutely continuous and the integral of the $L^\infty$ class $F$ is well defined.

Suppose $F_j,F$ are uniformly bounded and $F_j\to F$ in $L^1(\mathrm{d}V)$. If $\nu_j\in\mathcal A_\varphi(F_j)$, their densities are uniformly integrable. Every weak $L^1$ cluster point belongs to $\mathcal A_\varphi(F)$. If the latter is a singleton, the full sequence converges to it weakly in $L^1(\mathrm{d}V)$.
\end{theorem}
\begin{proof}
First, if $H_\varphi(\nu)<\infty$, the measure $\nu$ is absolutely continuous by \cref{thm:normalization-dual-compactness}. A bounded measurable $G$ admits uniformly bounded continuous approximations $G_k\to G$ in $L^1(\mathrm{d}V)$. Absolute continuity and the common bound imply $\int_XG_k\,\mathrm{d}\nu\to\int_XG\,\mathrm{d}\nu$. Together with \cref{thm:normalization-bounded-lone}, this proves
\begin{equation}\label{eq:dual-measurable-tests}
 H_\varphi(\nu)
 =\sup_{G\in L^\infty(\mathrm{d}V)}
       \left(B_\varphi(G)+\int_XG\,\mathrm{d}\nu\right)
\end{equation}
for such $\nu$. Values of measurable functions are understood modulo $\mathrm{d}V$-null sets. The finite-dual measures are absolutely continuous, so their integrals are well defined on these equivalence classes.

Choose continuous $F_j$ with $|F_j|\leq M=\|F\|_\infty$ and $F_j\to F$ in $L^1(\mathrm{d}V)$, and choose $\nu_j\in\mathcal A_\varphi(F_j)$. The preceding duality and Lipschitz bound give
\[
 H_\varphi(\nu_j)
 =B_\varphi(F_j)+\int_XF_j\,\mathrm{d}\nu_j
 \leq B_\varphi(0)+2M.
\]
Thus a subsequence converges weakly in $L^1$ to a probability measure $\nu$. Uniform absolute continuity gives
\[
 \int_X|F_j-F|\,\mathrm{d}\nu_j\longrightarrow0.
\]
Here again one splits at a fixed small positive threshold; the exceptional sets have $\mathrm{d}V$-measure tending to zero. Since $F$ is bounded measurable, weak $L^1$ convergence applies to its integral. Lower semicontinuity and continuity of $B_\varphi$ yield
\[
 H_\varphi(\nu)
 \leq\liminf_j H_\varphi(\nu_j)
 =B_\varphi(F)+\int_XF\,\mathrm{d}\nu.
\]
The reverse inequality follows from \eqref{eq:dual-measurable-tests}. Equality and that same supremum formula prove \eqref{eq:measurable-adjoint-support} and the minimum formula.

For continuous $F$, every measure in the old supporting set has finite dual value by \cref{prop:normalization-dual}, so \eqref{eq:dual-measurable-tests} extends all its inequalities to bounded measurable tests. This proves agreement of the definitions. Conversely, an absolutely continuous probability satisfying \eqref{eq:measurable-adjoint-support} has finite dual value and achieves the minimum. Hence the characterization is exact.

For the stability assertion, begin with the stated $F_j,\nu_j$ rather than with an approximating construction. The common uniform bound again places all $\nu_j$ in a single finite sublevel of $H_\varphi$. The same compactness and limiting argument show that every weak $L^1$ cluster point achieves the minimum for $F$. If the minimizing set is a singleton, subsequence compactness forces convergence of the full sequence to that measure. No strong $L^1$ convergence is asserted.
\end{proof}

The dual set therefore extends the smooth Gauduchon adjoint to bounded logarithmic densities, with a precise form of stability. Its possible failure to be a singleton remains a separate question.

\section{Entropy and the closed case}

We first recall a measure-theoretic identity used in both the closed and nonclosed settings.

For a probability measure $\nu=h\,\mathrm{d}V$, define
\[
 \operatorname{Ent}_{\mathrm{d}V}(\nu)=\int_Xh\log h\,\mathrm{d}V,
 \qquad 0\log0=0,
\]
and set its entropy to $+\infty$ if it is not absolutely continuous. This entropy may be negative when $\mathrm{d}V(X)\ne1$, but its negative part is finite since $h\log h\geq-\mathrm{e}^{-1}$.

\begin{lemma}\label{lem:entropy-variational}
For every $\nu\in\mathcal P(X)$,
\[
 \operatorname{Ent}_{\mathrm{d}V}(\nu)
 =\sup_{F\in C(X)}\left(\int_XF\,\mathrm{d}\nu
                     -\log\int_X\mathrm{e}^F\,\mathrm{d}V\right).
\]
\end{lemma}
\begin{proof}
If $\nu$ is not absolutely continuous, choose a compact set $K$ with $\mathrm{d}V(K)=0$ and $\nu(K)>0$. For $t>0$, regularity gives a continuous cutoff $0\leq\chi_t\leq1$, equal to one on $K$, whose support has volume at most $\mathrm{e}^{-2t}$. For $F=t\chi_t$, the expression in the supremum is at least
\[
 t\nu(K)-\log\bigl(\mathrm{d}V(X)+\mathrm{e}^{-t}\bigr),
\]
which tends to infinity with $t$.

Suppose now that $\nu=h\,\mathrm{d}V$. If its entropy is finite, Jensen's inequality for the logarithm gives
\[
 \int_XF\,\mathrm{d}\nu-\operatorname{Ent}_{\mathrm{d}V}(\nu)
 \leq\log\int_{\{h>0\}}\mathrm{e}^F\,\mathrm{d}V
 \leq\log\int_X\mathrm{e}^F\,\mathrm{d}V.
\]
For infinite entropy this upper bound is automatic. For the opposite inequality, set
\[
 F_k=\log\bigl(\max\{k^{-1},\min\{h,k\}\}\bigr),\qquad k\geq1.
\]
Then $\int_X\mathrm{e}^{F_k}\,\mathrm{d}V\to1$ and $\int_XF_k\,\mathrm{d}\nu\to\operatorname{Ent}_{\mathrm{d}V}(\nu)$, with the value $+\infty$ allowed: the positive parts of the integrands $hF_k$ increase to that of $h\log h$, and their negative parts are dominated by the integrable negative part of $h\log h$. For each fixed $k$, approximate $F_k$ in $L^1(\mathrm{d}V+\nu)$ by continuous functions with the same bounds; both integrals, including the exponential one, converge to those of $F_k$. Hence the supremum over continuous functions is at least the value of the expression at $F_k$, and letting $k\to\infty$ completes the proof.
\end{proof}


If $\theta$ is K\"ahler, the ordinary type $\varphi$ has a fixed non-pluripolar mass $m_\varphi>0$ under the standing hypotheses, by the closed-background theory of \cite{SGPT}. Integrating the equation gives
\begin{equation}\label{eq:normalization-closed-formula}
 B_\varphi(F)=\log m_\varphi-\log\int_X\mathrm{e}^F\,\mathrm{d}V,
 \qquad
 \mathcal A_\varphi(F)
 =\left\{\frac{\mathrm{e}^F\,\mathrm{d}V}{\int_X\mathrm{e}^F\,\mathrm{d}V}\right\}.
\end{equation}
The second identity follows by differentiating in bounded directions and testing the supporting inequality with both signs of the parameter.

\begin{proposition}\label{prop:normalization-dual-entropy}
In the closed case,
\[
 H_\varphi(\nu)=\log m_\varphi+\operatorname{Ent}_{\mathrm{d}V}(\nu).
\]
Thus the minimum formula becomes the usual entropy variational formula for the logarithmic integral in \eqref{eq:normalization-closed-formula}.
\end{proposition}
\begin{proof}
Substitute \eqref{eq:normalization-closed-formula} into the defining supremum and apply \cref{lem:entropy-variational}.
\end{proof}

The reference volume is part of this dual functional. If it is replaced by $\mathrm{d}V'=\mathrm{e}^a \mathrm{d}V$ for a continuous real $a$, then
\[
 B'_\varphi(F)=B_\varphi(F+a),\qquad
 H'_\varphi(\nu)=H_\varphi(\nu)-\int_Xa\,\mathrm{d}\nu.
\]
The first identity follows from the equation; the second follows by the change of variable $G=F+a$ in the defining supremum. The displayed identity for the dual is valid with infinite values as well. In a nonclosed background, $H_\varphi$ is still defined by the same supremum, but total mass need not be fixed on an ordinary type, so the entropy computation does not determine it.

\section{Entropy bounds}

Let $X$ be a connected compact complex manifold of dimension $n\geq1$, let $\omega$ be a Hermitian form, and let $\mathrm{d}V$ be a smooth positive probability volume form on $X$. Use the scalar $B_0$, supporting set $\mathcal A_0$ and dual $H_0$ from \cref{sec:normalization-scalar,sec:normalization-supports,sec:normalization-dual}, with background $\theta=\omega$ and reference $\varphi=0$. Write $\operatorname{Ent}_{\mathrm{d}V}$ for entropy as defined in \cref{sec:normalization-entropy}.

For a Hermitian background and reference zero, the dual functional need not equal entropy plus a constant. Its finite domain is nevertheless the finite-entropy domain.

\begin{proposition}\label{prop:hermitian-dual-entropy-bounds}
Let $\theta=\omega$ be Hermitian, let $\varphi=0$, and normalize $\mathrm{d}V(X)=1$. For every $q>1$ there is a finite real constant $a_q$, and there is a finite real constant $A$, such that every Borel probability measure $\nu$ satisfies
\[
 a_q+\frac1q\operatorname{Ent}_{\mathrm{d}V}(\nu)
 \leq H_0(\nu)
 \leq A+n\operatorname{Ent}_{\mathrm{d}V}(\nu).
\]
The inequalities include the value $+\infty$. Thus $H_0(\nu)<\infty$ if and only if $\nu$ has finite entropy. A family of probabilities has uniformly bounded $H_0$ if and only if its entropies are uniformly bounded. Neither mass property is assumed.
\end{proposition}
\begin{proof}
In dimension one the closed formula \eqref{eq:normalization-closed-formula} gives $H_0=\log\int_X\omega+\operatorname{Ent}_{\mathrm{d}V}$, so the assertions follow. Suppose $n\geq2$.

For $G\in C(X)$, apply \cref{thm:lp-potential-estimate} to a bounded solution with density $\mathrm{e}^{B_0(G)+G}$. It gives a constant $c_q>0$, independent of $G$, such that
\[
 B_0(G)\geq\log c_q-\frac1q\log\int_X\mathrm{e}^{qG}\,\mathrm{d}V.
\]
For the other bound, fix a smooth Gauduchon metric $\gamma$, rescaled so that $\int_X\omega\wedge\gamma^{n-1}=1$, and write
\[
 k_\gamma=\left(\frac{\gamma^n}{\mathrm{d}V}\right)^{(n-1)/n},
 \qquad m=\min_X k_\gamma>0.
\]
For smooth $G$, the smooth equation, matrix arithmetic--geometric mean and Stokes give
\[
 \mathrm{e}^{B_0(G)/n}\int_X\mathrm{e}^{G/n}k_\gamma\,\mathrm{d}V\leq1.
\]
Uniform approximation and continuity of $B_0$ extend this inequality to continuous $G$. Consequently
\[
 B_0(G)\leq-n\log m-n\log\int_X\mathrm{e}^{G/n}\,\mathrm{d}V.
\]

Take the defining suprema in the two scalar bounds and apply \cref{lem:entropy-variational}, with $F=qG$ and $F=G/n$, respectively. This proves both inequalities for every probability measure, including infinite values, with $a_q=\log c_q$ and $A=-n\log m$. Fix $q=2$ for the finite-domain and bounded-family assertions.
\end{proof}

In particular the compact dual sublevels of \cref{thm:normalization-dual-compactness} have uniformly bounded entropy in this Hermitian zero-reference regime. The estimate does not identify $H_0$ with entropy or give a fixed $L^p$ bound for the densities, with $p>1$.

\subsection{The mass bounds as entropy differences}

In a closed background the difference between $H_0$ and the entropy is constant. In general it measures how much the total Monge--Amp\`ere mass of bounded potentials varies: its infimum and supremum recover the lower and upper masses $v_-(\omega)$ and $v_+(\omega)$ of \cref{def:bmp}.

Let $\mathcal P_{\mathrm{ent}}(X)$ denote the set of probability measures of finite entropy relative to the fixed smooth probability volume $\mathrm{d}V$. By \cref{prop:hermitian-dual-entropy-bounds}, the difference $H_0(\nu)-\operatorname{Ent}_{\mathrm{d}V}(\nu)$ is a finite real number for every $\nu$ in this set.

\Needspace{13\baselineskip}
\begin{proposition}\label{prop:entropy-mass-extrema}\index{Monge--Ampere mass@Monge--Amp\`ere mass!entropy characterization}
For every Hermitian background $\omega$,
\[
 \begin{split}
 \inf_{\nu\in\mathcal P_{\mathrm{ent}}(X)}
      \bigl(H_0(\nu)-\operatorname{Ent}_{\mathrm{d}V}(\nu)\bigr)
       &=\log v_-(\omega),\\
 \sup_{\nu\in\mathcal P_{\mathrm{ent}}(X)}
      \bigl(H_0(\nu)-\operatorname{Ent}_{\mathrm{d}V}(\nu)\bigr)
       &=\log v_+(\omega).
 \end{split}
\]
The same extrema are obtained over smooth strictly positive probability measures. Here $\log0=-\infty$ and $\log(+\infty)=+\infty$. Thus the difference is bounded above if and only if $X$ has the bounded mass property, and bounded below if and only if $\omega$ has positive lower mass.
\end{proposition}
\begin{proof}
Write $Z_G=\int_X\mathrm{e}^G\,\mathrm{d}V$. A bounded solution for $G\in C(X)$ has mass $\mathrm{e}^{B_0(G)}Z_G$. If $v_-(\omega)>0$, the lower mass bound gives
\[
 B_0(G)\geq\log v_-(\omega)-\log Z_G.
\]
Taking the defining supremum for $H_0$ and using \cref{lem:entropy-variational} gives $H_0-\operatorname{Ent}_{\mathrm{d}V}\geq\log v_-(\omega)$ on $\mathcal P_{\mathrm{ent}}(X)$. In the same way, if $v_+(\omega)<\infty$, the upper mass bound gives $H_0-\operatorname{Ent}_{\mathrm{d}V}\leq\log v_+(\omega)$. These bounds hold trivially when $v_-(\omega)=0$ or $v_+(\omega)=+\infty$.

To prove sharpness, first take a smooth real $G$ and its Gibbs probability $\nu_G=Z_G^{-1}\mathrm{e}^G\,\mathrm{d}V$, which has finite entropy. Testing the supremum defining $H_0$ with $G$ gives
\[
 H_0(\nu_G)-\operatorname{Ent}_{\mathrm{d}V}(\nu_G)
       \geq B_0(G)+\log Z_G.
\]
In the other direction, let $\sigma_G$ be the smooth strictly positive Gauduchon adjoint probability of \cref{thm:normalization-gauduchon-recovery}. It lies in $\mathcal A_0(G)$, so its dual value is $B_0(G)+\int_XG\,\mathrm{d}\sigma_G$, and \cref{lem:entropy-variational} yields
\[
 H_0(\sigma_G)-\operatorname{Ent}_{\mathrm{d}V}(\sigma_G)
       \leq B_0(G)+\log Z_G.
\]
As $G$ runs over smooth real functions, the masses $\mathrm{e}^{B_0(G)}Z_G$ run over the masses $\int_X(\omega+\ddc u)^n$ of all smooth $u$ with $\omega+\ddc u>0$, since such a $u$ solves the equation for $G=\log((\omega+\ddc u)^n/\mathrm{d}V)$. By \cref{thm:bounded-exact-regularization}, the mass of every bounded $\omega$-plurisubharmonic function is a limit of such masses. Hence the infimum and supremum of $\mathrm{e}^{B_0(G)}Z_G$ over smooth $G$ are $v_-(\omega)$ and $v_+(\omega)$. Since $\nu_G$ and $\sigma_G$ are smooth and strictly positive, the last two inequalities prove sharpness on this smaller domain, including the cases $v_-(\omega)=0$ and $v_+(\omega)=+\infty$. The two equivalences follow.
\end{proof}

The difference does not depend on the smooth positive reference volume. For smooth real $a$, replacing $\mathrm{d}V$ by $\mathrm{d}V'=\mathrm{e}^a\mathrm{d}V$ changes both $H_0(\nu)$ and $\operatorname{Ent}(\nu)$ by $-\int_Xa\,\mathrm{d}\nu$, by the transformation formula preceding \cref{prop:hermitian-dual-entropy-bounds}, and finiteness of the entropy is unchanged because $a$ is bounded. Similarly, if $h$ is a smooth real function and $\omega+\ddc h$ is also Hermitian, transporting bounded solutions by $u\mapsto u-h$ preserves $B_0$, and hence this difference. Finally, the difference is constant on $\mathcal P_{\mathrm{ent}}(X)$ exactly when $v_-(\omega)=v_+(\omega)$, that is, when all bounded $\omega$-plurisubharmonic functions have the same mass.


\subsection{Entropy for a prescribed ordinary type}

Return to a semipositive big $\theta$ and a model reference $\varphi$ satisfying \eqref{eq:als-model-condition}. Fix a corresponding $\varepsilon\in(0,1)$. The unused part of the background supplies a uniform lower density bound; comparison in a larger Hermitian background supplies the upper bound. Together they extend the entropy estimates to the prescribed type.

\begin{corollary}\label{cor:prescribed-dual-entropy}
For every $q>1$ there is a finite real $a_q$, and there is a finite real $A$, such that
\[
 a_q+\frac1q\operatorname{Ent}_{\mathrm{d}V}(\nu)
 \leq H_\varphi(\nu)
 \leq A+n\operatorname{Ent}_{\mathrm{d}V}(\nu)
\]
for every Borel probability $\nu$, including infinite values. Thus finite dual value and finite entropy have the same domain, and uniformly bounded dual values are equivalent to uniformly bounded entropies. For fixed $\theta$ and $\varepsilon$ in \eqref{eq:als-model-condition}, the constants can be chosen independently of the model reference $\varphi$.
\end{corollary}
\begin{proof}
For $G\in C(X)$, let $u\sim\varphi$ solve the equation with density $c\mathrm{e}^G$, where $c=\mathrm{e}^{B_\varphi(G)}$. Applying \cref{lem:relative-density-lower-bound} gives $c\|\mathrm{e}^G\|_q\geq b_q$, hence
\[
 B_\varphi(G)\geq\log b_q-\frac1q
                         \log\int_X\mathrm{e}^{qG}\,\mathrm{d}V.
\]

For the upper bound choose a Hermitian form $\Omega\geq\theta$, and let $B_{\Omega,0}$ denote its bounded-type normalization. A bounded solution $v$ for $G$ in this background has a genuine smaller-background envelope. Moreover $u\in\PSH(X,\Omega)$, $u\preceq v$, and positivity gives
\[
 \mathrm{MA}^{\mathrm{ext}}_\Omega(u)
       \geq\mathrm{MA}^{\mathrm{ext}}_\theta(u)
       =c\mathrm{e}^G\,\mathrm{d}V.
\]
If $c>\mathrm{e}^{B_{\Omega,0}(G)}$, the same strict domination argument, now applied to $v$ and $u+t$ in background $\Omega$, is impossible. Hence $B_\varphi\leq B_{\Omega,0}$ and $H_\varphi\leq H_{\Omega,0}$. \Cref{prop:hermitian-dual-entropy-bounds} gives the required upper bound.

Taking the defining supremum in the lower scalar estimate and applying \cref{lem:entropy-variational} with $F=qG$ gives the lower bound for every probability, including infinite values, with $a_q=\log b_q$. The finite-domain and bounded-family assertions follow by taking $q=2$. The lower-density lemma is uniform in the reference for fixed $\varepsilon$, and the upper bound uses only $\Omega$. This proves the final assertion.
\end{proof}

For example, let the model references $\varphi_j$ satisfy the same smaller-background condition with fixed $\varepsilon$, let $G_j\in L^\infty(\mathrm{d}V)$ satisfy $\|G_j\|_\infty\leq M$ for a fixed constant $M$, and choose $\nu_j\in\mathcal A_{\varphi_j}(G_j)$. Choose a Hermitian form $\Omega\geq\theta$ and use its zero-reference functionals $B_{\Omega,0},H_{\Omega,0}$. Then
\[
 H_{\varphi_j}(\nu_j)
 =B_{\varphi_j}(G_j)+\int_XG_j\,\mathrm{d}\nu_j
 \leq B_{\Omega,0}(0)+2M.
\]
The uniform lower entropy estimate bounds their entropies. The Hermitian upper estimate places all $\nu_j$ in a single finite sublevel of $H_{\Omega,0}$. Hence \cref{thm:normalization-dual-compactness} gives weak compactness and uniform integrability of their densities. Every weak limit has finite entropy. This gives compactness as the reference varies; identifying the limit as a supporting probability for a limiting reference requires further information.




\section{Changing the prescribed singularities}

Keep the semipositive big background $\theta$ fixed, and use $B_0,H_0$ for its zero-reference functionals. A common smaller-background condition controls the change of normalization uniformly over all continuous logarithmic densities.

\begin{proposition}\label{prop:reference-normalization-comparison}
Let $\varphi,\psi$ satisfy \eqref{eq:normalization-reference}. If $\varphi\preceq\psi$, then $B_\varphi\leq B_\psi$ and $H_\varphi\leq H_\psi$. Moreover, if $P_{\varepsilon\theta}(\varphi)$ is genuine for $0<\varepsilon<1$, then
\[
 B_0(G)+n\log(1-\varepsilon)
       \leq B_\varphi(G)\leq B_0(G)
       \qquad(G\in C(X)),
\]
and
\[
 H_0(\nu)+n\log(1-\varepsilon)
       \leq H_\varphi(\nu)\leq H_0(\nu)
       \qquad(\nu\in\mathcal P(X)).
\]
The dual inequalities include infinite values.
\end{proposition}
\begin{proof}
Fix $G$, and let $u_\varphi,u_\psi$ solve the two ordinary-type equations, with scalars $c_\varphi,c_\psi$. Then $u_\varphi\preceq u_\psi$. If $c_\varphi>c_\psi$, singular strict domination applied to $u_\psi$ and $u_\varphi+t$ would give $u_\psi\geq u_\varphi+t$ for every real $t$, a contradiction at a common finite point. This proves the monotonicity of $B$, and its defining supremum gives the monotonicity of $H$. Since $\varphi\preceq0$, it also proves both upper bounds.

For the lower bound, let $v$ be the bounded solution for $G$, with scalar $c_0$, and put
\[
 h=P_{\varepsilon\theta}(\varphi)+(1-\varepsilon)v.
\]
The two currents in this sum have backgrounds $\varepsilon\theta$ and $(1-\varepsilon)\theta$, respectively. Positivity and homogeneity give
\[
 \mathrm{MA}^{\mathrm{ext}}_\theta(h)
       \geq(1-\varepsilon)^n\mathrm{MA}_\theta(v)
       =(1-\varepsilon)^nc_0\mathrm{e}^G\,\mathrm{d}V.
\]
As in the lower-density lemma, the inequality is justified first on finite levels and then off the pole set of $P_{\varepsilon\theta}(\varphi)$, which has zero smooth volume. Also $h\preceq u_\varphi$. If $c_\varphi<(1-\varepsilon)^nc_0$, strict domination would force $u_\varphi\geq h+t$ for every $t$. Thus $c_\varphi\geq(1-\varepsilon)^nc_0$. Taking logarithms proves the scalar bound; taking the defining suprema proves the dual bound.
\end{proof}

\begin{corollary}\label{cor:small-reference-adjoint-limits}
Let the model references $\varphi_j$ satisfy $P_{\varepsilon_j\theta}(\varphi_j)\not\equiv-\infty$, where $0<\varepsilon_j<1$ and $\varepsilon_j\to0$. Then $B_{\varphi_j}\to B_0$ uniformly on $C(X)$. If $G_j\to G$ uniformly and $\nu_j\in\mathcal A_{\varphi_j}(G_j)$, every weak cluster point belongs to $\mathcal A_0(G)$. The densities of $\nu_j$ are uniformly integrable, and weak convergence of a subsequence also holds weakly in $L^1(\mathrm{d}V)$. If $\mathcal A_0(G)$ is a singleton, the full sequence converges to it in these two senses.
\end{corollary}
\begin{proof}
Put $\delta_j=-n\log(1-\varepsilon_j)\to0$. The proposition gives $|B_{\varphi_j}-B_0|\leq\delta_j$ on all of $C(X)$. Its dual inequality and the supporting identity give
\[
 0\leq H_0(\nu_j)-\int_XG_j\,\mathrm{d}\nu_j-B_0(G_j)
       \leq\delta_j.
\]
Choose $M$ with $\|G_j\|_\infty\leq M$. This implies $H_0(\nu_j)\leq B_0(0)+2M+\delta_j$, so all $\nu_j$ lie in one compact dual sublevel. If a subsequence converges weakly to $\nu$, lower semicontinuity, uniform convergence of $G_j$ and continuity of $B_0$ give
\[
 H_0(\nu)-\int_XG\,\mathrm{d}\nu-B_0(G)\leq0.
\]
The dual minimum formula gives the opposite inequality and identifies equality with $\nu\in\mathcal A_0(G)$. The common-sublevel theorem \cref{thm:normalization-dual-compactness} also gives uniform integrability and weak $L^1$ convergence. Compactness and uniqueness of the limiting supporting probability give the last assertion.
\end{proof}


\Needspace{8\baselineskip}
\subsection{Generic uniqueness}
\label{sec:generic-adjoints}

On every affine line in $C(X)$, the concave function $B_\varphi$ is differentiable outside a countable set, so the gap between its two directional derivatives along the line vanishes on a dense subset. Upper semicontinuity of this gap and separability of $C(X)$ turn the one-dimensional fact into uniqueness of the supporting probability on a dense $G_\delta$ subset of $C(X)$.

\Needspace{12\baselineskip}
\begin{theorem}\label{thm:generic-adjoint-uniqueness}\label{cor:generic-adjoint-approximation}\index{Adjoint measure!generic uniqueness}
The set
\[
 \mathscr U_\varphi
   =\{F\in C(X,\mathbb R):\mathcal A_\varphi(F)
                                  \text{ is a singleton}\}
\]
is a dense $G_\delta$ subset of $C(X)$ with the uniform norm. At every $F\in\mathscr U_\varphi$, the functional $B_\varphi$ is G\^ateaux differentiable, and its derivative is minus the integral against the unique supporting probability. If $F_j\to F$ uniformly and $\nu_j\in\mathcal A_\varphi(F_j)$, then $\nu_j$ converges weakly to that probability.
\end{theorem}
\Needspace{5\baselineskip}
\begin{proof}
For $h\in C(X,\mathbb R)$ define the width
\[
 w_h(F)=\max_{\nu\in\mathcal A_\varphi(F)}\int_Xh\,\mathrm{d}\nu
         -\min_{\nu\in\mathcal A_\varphi(F)}\int_Xh\,\mathrm{d}\nu\geq0.
\]
The extrema exist by weak compactness. We claim that the width is upper semicontinuous in $F$. Indeed, let $F_j\to F$ and choose maximizing and minimizing probabilities at each $F_j$. Along a subsequence realizing the upper limit of the widths, compactness of the space of probability measures gives weak limits of both choices. By \cref{prop:weak-adjoint-stability}, both limits belong to $\mathcal A_\varphi(F)$, so the difference of their integrals against $h$ is at most $w_h(F)$.

Consequently, for every $\epsilon>0$ the set $\{F:w_h(F)<\epsilon\}$ is open. It is also dense. If $h=0$ there is nothing to prove. Otherwise fix $F$ and an open ball about it. The function
\[
 b(t)=B_\varphi(F+th)
\]
is finite and concave on the real line, so it is differentiable outside an at most countable set: its one-sided derivatives are finite and monotone, and the nonempty open intervals between unequal one-sided derivatives at distinct points are disjoint, so each contains a different rational number. Choose a point of differentiability $t$ so close to zero that $F+th$ lies in the given ball. The derivative formulas of \cref{prop:normalization-directional} give $w_h(F+th)=0$, which proves density.

The compact manifold $X$ is metrizable, so $C(X)$ is separable. Choose a countable uniformly dense family $(h_j)$ in $C(X)$. By the Baire category theorem in the complete metric space $C(X)$, the set
\[
 \bigcap_{j,k\geq1}\{F:w_{h_j}(F)<1/k\}
\]
is a dense $G_\delta$. At every point of this intersection, all supporting probabilities have the same integral against every $h_j$. Since they have unit mass, uniform approximation gives equal integrals against every continuous function, and the Riesz representation theorem shows that the probabilities coincide. Conversely, a singleton has zero width in every direction, so the intersection is exactly $\mathscr U_\varphi$. The formula for the derivative follows from \cref{prop:normalization-directional}, and the approximation assertion from \cref{prop:weak-adjoint-stability}.
\end{proof}


\section{Selection by a vanishing exponential parameter}
\label{sec:vanishing-parameter-selection}

The equation with exponent zero can have several normalized solutions. Nevertheless, the solutions of the equations with positive exponent select one of them in the limit: the scalar functional provides a family of integral normalizations, and the selected potential is the greatest subsolution of the critical equation satisfying all of them. To prove that it is the greatest, we use supporting probabilities of a second functional, obtained by varying the forcing term in the subsolution inequality.

In this section, $\theta$ is an arbitrary smooth big real $(1,1)$-form on $X$, not necessarily semipositive. Fix a strict analytic barrier $\rho\leq0$ as in \cref{sec:big-minimal-comparison}, and put $\Omega=X\setminus\{\rho=-\infty\}$. A potential has \emph{minimal type} if it has the ordinary singularity type of the envelope $V_\theta=P_\theta(0)$, which has minimal singularities (\cref{sec:obstacle-envelopes}). Fix a nonzero $0\leq f\in L^p(\mathrm{d}V)$, with $1<p<\infty$. For a nonzero nonnegative density $g$ in some $L^q$, $q>1$, let $c(g)>0$ be the unique constant of \cref{thm:big-minimal-normalized} for which the equation $\mathrm{MA}_\theta(v)=c(g)g\,\mathrm{d}V$ has a solution $v$ of minimal type. We call $c(g)$ the \emph{critical scalar} of $g$ and such solutions \emph{critical potentials} for $g$; they are continuous on $\Omega$ by \cref{thm:big-minimal-holder}. Here $\mathrm{MA}_\theta$ is the extended measure $\mathrm{MA}^{\mathrm{ext}}_\theta$. On $\Omega$ it is the locally bounded product for every potential of minimal type, and it gives zero mass to the analytic set $X\setminus\Omega$.

Put $c=c(f)$ and define
\begin{equation}\label{eq:selection-probabilities}
 \begin{aligned}
 \kappa(r)&=\log c-\log c(\mathrm{e}^r f),\\
 \mathcal P_f&=\left\{\nu\in\mathcal P(X):
       \int_X r\,\mathrm{d}\nu\leq\kappa(r)
       \text{ for every }r\in C(X)\right\}.
 \end{aligned}
\end{equation}
When $\theta$ is semipositive and $f=\mathrm{e}^F$ with $F\in C(X)$, we have $\kappa(r)=B_0(F)-B_0(F+r)$, and $\mathcal P_f$ is the set $\mathcal A_0(F)$ of supporting probabilities. The definition also applies when $f$ vanishes on a set of positive measure and when $\theta$ is not semipositive.

\begin{theorem}\label{thm:vanishing-exponential-selection}\index{Monge--Ampere equation@Monge--Amp\`ere equation!vanishing-exponent selection}
For $\lambda>0$, let $u_\lambda$ be the solution of minimal type of
\[
 \mathrm{MA}_\theta(u_\lambda)
       =\mathrm{e}^{\lambda u_\lambda}f\,\mathrm{d}V
\]
given by \cref{thm:big-minimal-exponential} and unique by \cref{cor:big-barrier-orders}. The set $\mathcal P_f$ is nonempty and weakly compact, and every potential of minimal type is integrable against all its measures. As $\lambda\downarrow0$,
\[
 w_\lambda=u_\lambda-\lambda^{-1}\log c
\]
converges locally uniformly on $\Omega$ and in $L^1(X)$ to a potential $w$ of minimal type, continuous on $\Omega$. It is the greatest among the potentials $s$ of minimal type, continuous on $\Omega$, satisfying
\begin{equation}\label{eq:selection-class}
 \mathrm{MA}_\theta(s)\geq cf\,\mathrm{d}V,\qquad
 \int_Xs\,\mathrm{d}\nu\leq0\quad(\nu\in\mathcal P_f).
\end{equation}
Moreover,
\[
 \mathrm{MA}_\theta(w)=cf\,\mathrm{d}V,\qquad
 \max_{\nu\in\mathcal P_f}\int_Xw\,\mathrm{d}\nu=0.
\]
Consequently, $u_\lambda-\sup_Xu_\lambda$ converges to $w-\sup_Xw$ in the same two topologies. For a Hermitian $\theta$, one may take $\Omega=X$, and both convergences of potentials are uniform.
\end{theorem}

\subsection{A norm defined by the critical scalar}

Mixed products turn the comparison of critical scalars into a triangle inequality for $n$th roots of densities, and the subsolution estimate associated with the strict barrier gives continuity in the corresponding Lebesgue norm.

Fix $1<q<p$ and set $\ell=nq$. On real $L^\ell(\mathrm{d}V)$ define
\[
 N(a)=
 \begin{cases}
 c(|a|^n)^{-1/n},&a\ne0,\\
 0,&a=0.
 \end{cases}
\]

\Needspace{12\baselineskip}
\begin{proposition}\label{prop:critical-scalar-norm}
The function $N$ is a lattice norm on $L^\ell(\mathrm{d}V)$, that is, a norm satisfying $N(a)\leq N(b)$ whenever $|a|\leq|b|$. There is $C_q<\infty$, depending only on $X,\theta,\mathrm{d}V$ and $q$, such that
\begin{equation}\label{eq:critical-scalar-norm-bound}
 N(a)\leq C_q\|a\|_\ell,\qquad
 |N(a)-N(b)|\leq C_q\|a-b\|_\ell.
\end{equation}
Consequently, for nonzero $g,h\geq0$ in $L^q$,
\[
 |c(g)^{-1/n}-c(h)^{-1/n}|
       \leq C_q\|g-h\|_q^{1/n}.
\]
\end{proposition}
\begin{proof}
If a potential $v$ of minimal type satisfies $\mathrm{MA}_\theta(v)\geq d g\,\mathrm{d}V$, then $d\leq c(g)$, since otherwise comparison with a critical potential for $g$ contradicts part (1) of \cref{cor:big-barrier-orders}. Hence $g\leq h$ implies $c(g)\geq c(h)$. Uniqueness of the critical scalar also gives $c(tg)=c(g)/t$ for $t>0$. These facts prove that $N$ is monotone and absolutely homogeneous, and positivity follows from $0<c(g)<\infty$.

For the triangle inequality, let $a,b\geq0$ be nonzero, put $A=N(a)$ and $B=N(b)$, and choose critical potentials $v_a,v_b$ with measures $(a/A)^n\mathrm{d}V$ and $(b/B)^n\mathrm{d}V$. On every relatively compact subset of $\Omega$, \eqref{eq:interior-minkowski-background} gives
\[
 \mathrm{MA}_\theta\left(\frac{A v_a+B v_b}{A+B}\right)
       \geq\left(\frac{a+b}{A+B}\right)^n\mathrm{d}V,
\]
where both potentials are locally bounded. Exhausting $\Omega$ gives the inequality on $X$, since $X\setminus\Omega$ has zero smooth volume. The averaged potential has minimal type, so the first paragraph gives $N(a+b)\leq A+B$. The cases where $a$ or $b$ vanishes are trivial, and for real functions we use $|a+b|\leq|a|+|b|$ and monotonicity.

Next, apply \cref{lem:big-reference-subsolution} with $\eta=V_\theta$ to nonnegative densities of $L^q$ norm at most one. It gives a uniform $m>0$ and subsolutions of minimal type, hence $c(g)\geq m$ for each such nonzero density. For $a\ne0$, take $g=|a|^n/\|a\|_\ell^n$, which has $L^q$ norm one. Homogeneity gives $N(a)\leq m^{-1/n}\|a\|_\ell$, which is the first bound in \eqref{eq:critical-scalar-norm-bound}, and the reverse triangle inequality gives the second. Applying the second bound to $g^{1/n},h^{1/n}$ and using $|g^{1/n}-h^{1/n}|^n\leq|g-h|$ proves the last assertion.
\end{proof}

\subsection{Norming probabilities and integral bounds}

Set $a=f^{1/n}\in L^\ell(\mathrm{d}V)$ and $A=N(a)=c^{-1/n}$. The probabilities in \eqref{eq:selection-probabilities} come from positive linear functionals that norm $N$ at $a$, that is, functionals dominated by $N$ and equal to $N(a)$ at $a$. This representation will control integrals of unbounded potentials as well as bounded logarithmic perturbations.

With $\ell'=\ell/(\ell-1)$, define
\begin{equation}\label{eq:selection-norming-face}
 \mathcal H_a=\left\{\eta\in L^{\ell'}(\mathrm{d}V):
 \eta\geq0,\ \int_Xa\eta\,\mathrm{d}V=A,\
 \int_Xb\eta\,\mathrm{d}V\leq N(b)\text{ for every }b\in L^\ell\right\}.
\end{equation}

\Needspace{18\baselineskip}
\begin{proposition}\label{prop:selection-norming-probabilities}
The set $\mathcal H_a$ is nonempty and weakly compact in $L^{\ell'}$. The function $\kappa$ of \eqref{eq:selection-probabilities} is given by
\[
 \kappa(r)=n\log\frac{N(a\mathrm{e}^{r/n})}{A},
\]
and it is convex, increasing and $1$-Lipschitz for the uniform norm on $C(X)$, with $\kappa(0)=0$ and $\kappa(r+b)=\kappa(r)+b$ for constants $b$. Its subdifferential at zero is
\begin{equation}\label{eq:selection-face-representation}
 \mathcal P_f=\{\nu_\eta:\eta\in\mathcal H_a\},\qquad
 \mathrm{d}\nu_\eta=A^{-1}a\eta\,\mathrm{d}V.
\end{equation}
If $z$ is measurable and $az\in L^\ell$, then
\begin{equation}\label{eq:selection-weighted-moment}
 \sup_{\nu\in\mathcal P_f}\int_X|z|\,\mathrm{d}\nu
       \leq C_q A^{-1}\|az\|_\ell.
\end{equation}
If also $a\mathrm{e}^{tz/n}\in L^\ell$ and $N(a\mathrm{e}^{tz/n})=A$ for some $t>0$, then $\int_Xz\,\mathrm{d}\nu\leq0$ for every $\nu\in\mathcal P_f$.
\end{proposition}
\begin{proof}
The formula for $\kappa$ follows from $N(a\mathrm{e}^{r/n})=c(\mathrm{e}^rf)^{-1/n}$. The Hahn--Banach theorem gives a linear functional that norms $N$ at $a$. By \eqref{eq:critical-scalar-norm-bound} it is continuous on $L^\ell$, and hence is represented by some $\eta_0\in L^{\ell'}$. For $b\geq0$, testing on $\operatorname{sgn}(\eta_0)b$ gives
\[
 \int_Xb|\eta_0|\,\mathrm{d}V\leq N(b).
\]
Since $a\geq0$ and $\int_Xa\eta_0\,\mathrm{d}V=A=N(a)$, the same equality holds with $|\eta_0|$ in place of $\eta_0$. Monotonicity of $N$ extends the domination to real $b$, so $|\eta_0|\in\mathcal H_a$. The defining conditions of $\mathcal H_a$ are weakly closed and give $\|\eta\|_{\ell'}\leq C_q$, so reflexivity proves weak compactness.

For nonzero nonnegative $b_0,b_1$, the weighted arithmetic--geometric mean inequality applied to $b_j/N(b_j)$, together with monotonicity and the triangle inequality, gives
\[
 N(b_0^{1-t}b_1^t)\leq N(b_0)^{1-t}N(b_1)^t,\qquad 0\leq t\leq1.
\]
Taking $b_j=a\mathrm{e}^{r_j/n}$ proves convexity of $\kappa$. Monotonicity and the translation rule follow from the corresponding properties of $N$, and together they imply the Lipschitz bound.

\Needspace{5\baselineskip}
Every $\nu_\eta$ is a probability and satisfies
\[
 \frac{N(a\mathrm{e}^{r/n})}{A}
       \geq\int_X\mathrm{e}^{r/n}\,\mathrm{d}\nu_\eta.
\]
Jensen's inequality gives $\kappa(r)\geq\int_Xr\,\mathrm{d}\nu_\eta$, which proves one inclusion in \eqref{eq:selection-face-representation}. The map $\eta\mapsto\nu_\eta$ is weakly continuous, since $ar\in L^\ell$ for $r\in C(X)$.

Conversely, let $\nu\in\mathcal P_f$ and put $L(r)=\int_Xr\,\mathrm{d}\nu$. Differentiating the curve $t\mapsto a\mathrm{e}^{tr/n}$ in $L^\ell$ at $t=0$ gives
\[
 N(a\mathrm{e}^{tr/n})
       =N(a)+\frac{t}{n}N'_+(a;ar)+o(t),\qquad t\downarrow0.
\]
Here the one-sided derivative $N'_+(a;\cdot)$ of the continuous convex function $N$ exists, and the error estimate follows from its Lipschitz bound. The inequality $L\leq\kappa$, applied to $tr$ and $-tr$, therefore gives
\begin{equation}\label{eq:selection-face-extension}
 |A L(r)|\leq C_q\|ar\|_\ell,\qquad
 A L(r)\leq N'_+(a;ar)\leq N(a+ar)-A.
\end{equation}
Thus $T(ar)=A L(r)$ defines a continuous linear functional on the space of products $ar$, $r\in C(X)$. These products are dense in
\[
 S_a=\{b\in L^\ell:b=0\text{ almost everywhere on }\{a=0\}\}:
\]
indeed, it suffices to approximate $b/a$ in $L^\ell(a^\ell\mathrm{d}V)$ by continuous functions. If $ar\geq0$, then $ar=ar_+$ almost everywhere, and the first inequality in \eqref{eq:selection-face-extension} gives $L(r-r_+)=0$, so $T(ar)\geq0$. Approximating nonnegative elements of $S_a$ by products $ar$ with $r\geq0$ continuous shows that the continuous extension of $T$ to $S_a$ is positive.

Represent this extension by some $\eta\geq0$ in $L^{\ell'}$ vanishing on $\{a=0\}$. Since $T(a)=A$, extending the second inequality in \eqref{eq:selection-face-extension} to $S_a$ and applying it to $b-a$ gives $T(b)\leq N(b)$ for $b\in S_a$. For general $b\in L^\ell$,
\[
 \int_Xb\eta\,\mathrm{d}V
 \leq N(b\mathbf1_{\{a>0\}})\leq N(b).
\]
Thus $\eta\in\mathcal H_a$, and $\nu_\eta$ agrees with $\nu$ on all continuous functions. This proves the representation \eqref{eq:selection-face-representation}, and weak compactness of $\mathcal P_f$ follows from that of $\mathcal H_a$.

The bound \eqref{eq:selection-weighted-moment} follows from H\"older's inequality and $\|\eta\|_{\ell'}\leq C_q$. Under the last hypothesis, the norming inequality gives $\int_X\mathrm{e}^{tz/n}\,\mathrm{d}\nu\leq1$. Since $z$ is $\nu$-integrable by \eqref{eq:selection-weighted-moment}, Jensen's inequality implies $\int_Xz\,\mathrm{d}\nu\leq0$.
\end{proof}

These estimates apply uniformly to all $v\in\PSH(X,\theta)$ with $\sup_Xv=0$. Indeed, \cref{lem:uniform-exponential}, applied in a Hermitian background larger than $\theta$, gives uniform bounds for all finite moments of $|v|$, and H\"older's inequality gives
\begin{equation}\label{eq:selection-psh-moments}
 \int_X |v|^{nq} f^q\,\mathrm{d}V
 \leq \|f\|_p^q
       \left(\int_X|v|^{nqp/(p-q)}\,\mathrm{d}V\right)^{(p-q)/p}.
\end{equation}
Consequently, \eqref{eq:selection-weighted-moment} gives a uniform bound for $\int_X|v|\,\mathrm{d}\nu$ over all such $v$ and all $\nu\in\mathcal P_f$; constant translates of such $v$ are integrable as well. The representation also shows that these probabilities charge neither $\{f=0\}$ nor $X\setminus\Omega$.

\subsection{Bounds for the centered potentials}

We call the potentials $w_\lambda$ of \cref{thm:vanishing-exponential-selection} the \emph{centered potentials}. Choose a critical potential $\psi$ for $f$ with $\sup_X\psi=0$, so that $\mathrm{MA}_\theta(\psi)=cf\,\mathrm{d}V$. The integral normalizations prevent the centered potentials from drifting upward by constants.

\Needspace{10\baselineskip}
\begin{lemma}\label{lem:selection-centered-bounds}
There is $K<\infty$, independent of $0<\lambda\leq1$, such that
\begin{equation}\label{eq:selection-centered-bounds}
 \psi\leq w_\lambda,\qquad
 0\leq k_\lambda:=\sup_Xw_\lambda\leq K,\qquad
 \psi-K\leq v_\lambda:=w_\lambda-k_\lambda\leq V_\theta.
\end{equation}
Moreover $\int_Xw_\lambda\,\mathrm{d}\nu\leq0$ for every $\nu\in\mathcal P_f$. The family $w_\lambda$ is relatively compact for locally uniform convergence on $\Omega$, and every subsequence converging locally uniformly on $\Omega$ also converges in $L^1(X)$, to a potential of minimal type.
\end{lemma}
\begin{proof}
The centered potential satisfies $\mathrm{MA}_\theta(w_\lambda)=c\mathrm{e}^{\lambda w_\lambda}f\,\mathrm{d}V$, while $\psi\leq0$ gives $\mathrm{MA}_\theta(\psi)\geq c\mathrm{e}^{\lambda\psi}f\,\mathrm{d}V$. Part (2) of \cref{cor:big-barrier-orders} therefore gives $\psi\leq w_\lambda$, hence $k_\lambda\geq0$. For each fixed $\lambda$, the density $\mathrm{e}^{\lambda w_\lambda}f$ lies in $L^p$, and its critical scalar is $c$. Therefore
\[
 N(a\mathrm{e}^{\lambda w_\lambda/n})=A.
\]
The moment estimate above gives $aw_\lambda\in L^\ell$, so \cref{prop:selection-norming-probabilities}, with $z=w_\lambda$ and $t=\lambda$, gives the asserted integral inequalities. Fix one $\nu\in\mathcal P_f$. Then
\[
 k_\lambda\leq-\int_Xv_\lambda\,\mathrm{d}\nu
       \leq C_q A^{-1}\|a v_\lambda\|_\ell\leq K
\]
by \eqref{eq:selection-psh-moments}, with $K$ independent of $\lambda$. The normalization $\sup_Xv_\lambda=0$ gives $v_\lambda\leq V_\theta$, and subtracting $k_\lambda$ from the lower bound $\psi\leq w_\lambda$ proves \eqref{eq:selection-centered-bounds}.

The potentials $v_\lambda=u_\lambda-\sup_Xu_\lambda$ form the normalized exponential family of \cref{prop:big-exponential-clusters}, which gives subsequences converging locally uniformly on $\Omega$ and in $L^1(X)$ to potentials of minimal type. Extracting a convergent subsequence of $k_\lambda\in[0,K]$ gives the same compactness for $w_\lambda=v_\lambda+k_\lambda$. The bound $|w_\lambda|\leq|\psi|+K$ also gives $L^1$ convergence along every subsequence converging locally uniformly on $\Omega$, since $X\setminus\Omega$ has zero volume.
\end{proof}

\Needspace{4\baselineskip}
In particular,
\[
 \mathrm{e}^{\lambda\sup_Xu_\lambda}
       =c\mathrm{e}^{\lambda k_\lambda}\longrightarrow c.
\]
\Needspace{10\baselineskip}
Moreover, for $0<\lambda\leq1$ the densities of the equations satisfy $\|c\mathrm{e}^{\lambda w_\lambda}f\|_p\leq c\mathrm{e}^K\|f\|_p$. It matters that this bound includes the scalar factor: an $L^p$ bound on the data alone, after an arbitrary rescaling, does not bound the normalized potentials. For example, on $\mathbb P^1$, with $\omega_{\mathrm{FS}}=\ddc\log(1+|z|^2)$ and $0<\epsilon<1$, the function
\[
 \psi_\epsilon=\log\frac{|z|^2+\epsilon^2}{1+|z|^2}
\]
is smooth and $\omega_{\mathrm{FS}}$-plurisubharmonic, with supremum zero and infimum $\log\epsilon^2$. Write $\mathrm{MA}_{\omega_{\mathrm{FS}}}(\psi_\epsilon)=g_\epsilon\omega_{\mathrm{FS}}$, and take $L^p$ norms with respect to $\omega_{\mathrm{FS}}$ in this example. The data $g_\epsilon/\|g_\epsilon\|_p$ have norm one, and $\psi_\epsilon$ solves the equation with these data and scalar $\|g_\epsilon\|_p$; but the densities $g_\epsilon$ of the equations themselves are not bounded in $L^p$, and the potentials are not uniformly bounded.

\subsection{Scalar envelopes of subsolutions}

To compare a cluster potential with every member of the class \eqref{eq:selection-class}, fix a point $x\in\Omega$ and add a continuous forcing term $r$ to the subsolution inequality. Let $\mathcal S$ be the convex set of potentials of minimal type that are continuous on $\Omega$, and define
\begin{equation}\label{eq:selection-scalar-envelope}
 \mathcal F_{\lambda,x}(r)=
 \sup\{s(x):s\in\mathcal S,\
       \mathrm{MA}_\theta(s)\geq
              c\mathrm{e}^{\lambda s+r}f\,\mathrm{d}V\}.
\end{equation}

\begin{lemma}\label{lem:selection-forcing-probability}
The functional $\mathcal F_{\lambda,x}$ is finite, concave, decreasing and $\lambda^{-1}$-Lipschitz for the uniform norm on $C(X)$. It satisfies
\[
 \mathcal F_{\lambda,x}(0)=w_\lambda(x),\qquad
 \mathcal F_{\lambda,x}(r+b)
       =\mathcal F_{\lambda,x}(r)-b/\lambda.
\]
There is a probability measure $\nu_{\lambda,x}$ such that, for every $r\in C(X)$ and every $s$ admissible for $r$ in \eqref{eq:selection-scalar-envelope},
\begin{equation}\label{eq:selection-forcing-support}
 \lambda s(x)+\int_Xr\,\mathrm{d}\nu_{\lambda,x}
       \leq\lambda w_\lambda(x).
\end{equation}
\end{lemma}
\begin{proof}
By part (2) of \cref{cor:big-barrier-orders}, $w_\lambda$ is the greatest admissible potential for $r=0$. If $m=\min_Xr$ and $M=\max_Xr$, then $w_\lambda-M/\lambda$ is admissible for $r$, and translating any potential admissible for $r$ by $m/\lambda$ makes it admissible for $r=0$. Hence
\[
 w_\lambda(x)-M/\lambda\leq\mathcal F_{\lambda,x}(r)
       \leq w_\lambda(x)-m/\lambda.
\]
The mixed density inequality on $\Omega$ and the arithmetic--geometric mean inequality show that a convex combination of potentials admissible for $r_0$ and $r_1$ is admissible for the same convex combination of the forcing terms. Approximating the two suprema at $x$ proves concavity. Inclusion of the admissible classes proves monotonicity, and translation by constants proves the translation rule and the Lipschitz bound.

Now apply the Hahn--Banach construction in the proof of \cref{thm:weak-adjoint-probabilities} to $-\lambda\mathcal F_{\lambda,x}$ at zero. This functional is finite, continuous, convex and increasing, and it satisfies $-\lambda\mathcal F_{\lambda,x}(r+b)=-\lambda\mathcal F_{\lambda,x}(r)+b$ for constants $b$. The same tests with nonnegative functions and constants show that its subgradient is a probability measure $\nu_{\lambda,x}$. The resulting supporting inequality, combined with $s(x)\leq\mathcal F_{\lambda,x}(r)$, gives \eqref{eq:selection-forcing-support}; it is not necessary that the supremum be attained.
\end{proof}

\begin{proof}[of \cref{thm:vanishing-exponential-selection}]
By \cref{prop:selection-norming-probabilities} and the moment bound following it, $\mathcal P_f$ has the stated properties, and \cref{lem:selection-centered-bounds} gives precompactness of the centered potentials. It remains to identify every cluster potential.

\emph{Step 1: limits of the forcing probabilities.} Fix $x\in\Omega$. For $r\in C(X)$, choose a critical potential $\psi_r$ for $\mathrm{e}^r f$ with supremum zero. By the definition of $\kappa$,
\[
 \mathrm{MA}_\theta(\psi_r)
       =c\mathrm{e}^{-\kappa(r)+r}f\,\mathrm{d}V.
\]
Since $\psi_r\leq0$, the potential $\psi_r-\kappa(r)/\lambda$ is admissible in \eqref{eq:selection-scalar-envelope}. Thus
\begin{equation}\label{eq:selection-forcing-limit}
 \int_Xr\,\mathrm{d}\nu_{\lambda,x}
       \leq\kappa(r)+\lambda\bigl(w_\lambda(x)-\psi_r(x)\bigr).
\end{equation}
For fixed $x$ and $r$, the values $w_\lambda(x)$ are uniformly bounded by \cref{lem:selection-centered-bounds}, and $\psi_r(x)$ is finite. Hence every weak cluster point of $\nu_{\lambda,x}$ as $\lambda\downarrow0$ belongs to $\mathcal P_f$.

\emph{Step 2: a pointwise comparison.} Let $s$ satisfy \eqref{eq:selection-class}. The function $-\lambda s$ is lower semicontinuous and bounded below on $X$, with the value $+\infty$ allowed, so metric inf-convolution and truncation give continuous functions $r_j\uparrow-\lambda s$. Since $r_j\leq-\lambda s$, the potential $s$ is admissible for every $r_j$, and \eqref{eq:selection-forcing-support} bounds the integrals of the $r_j$ from above. Monotone convergence, after subtracting a common lower bound, shows that $s$ is integrable against $\nu_{\lambda,x}$ and gives
\begin{equation}\label{eq:selection-pointwise-comparison}
 w_\lambda(x)\geq s(x)-\int_Xs\,\mathrm{d}\nu_{\lambda,x}.
\end{equation}

\emph{Step 3: cluster potentials belong to the class \eqref{eq:selection-class}.} Take a sequence $\lambda_j\downarrow0$ with $w_{\lambda_j}\to z$ locally uniformly on $\Omega$ and in $L^1(X)$. The bounds \eqref{eq:selection-centered-bounds} show that $z$ has minimal type, and $z$ is continuous on $\Omega$. Since $c\mathrm{e}^{\lambda_j w_{\lambda_j}}f\to cf$ in $L^p$, by dominated convergence and the common bound $w_\lambda\leq K$, continuity of local bounded products gives $\mathrm{MA}_\theta(z)=cf\,\mathrm{d}V$ on $\Omega$, and then on $X$ because neither side charges $X\setminus\Omega$.

For each $\nu\in\mathcal P_f$, the bound $|w_\lambda|\leq|\psi|+K$ has a $\nu$-integrable right-hand side, and $\nu$ charges only $\Omega$. Dominated convergence in $\int_Xw_\lambda\,\mathrm{d}\nu\leq0$ therefore gives $\int_Xz\,\mathrm{d}\nu\leq0$. Hence $z$ belongs to the class \eqref{eq:selection-class}.

\emph{Step 4: $z$ is the greatest element of the class.} For any $s$ in the class and fixed $x\in\Omega$, pass to a further subsequence with $\nu_{\lambda_j,x}\rightharpoonup\nu$, so that $\nu\in\mathcal P_f$ by Step 1. The potential $s$ is upper semicontinuous and bounded above on $X$, so
\[
 \limsup_j\int_Xs\,\mathrm{d}\nu_{\lambda_j,x}
       \leq\int_Xs\,\mathrm{d}\nu\leq0.
\]
For an unbounded $s$, this follows from the portmanteau theorem applied to $\max(s,-M)$, followed by $M\to\infty$; the final integral is finite by \eqref{eq:selection-weighted-moment}. Passing to the limit in \eqref{eq:selection-pointwise-comparison} gives $z(x)\geq s(x)$. This holds for every $x\in\Omega$, and hence on $X$ by \cref{lem:qpsh-almost-everywhere-order}. Thus every cluster potential is the greatest element $w$ of the class, and precompactness gives convergence of the whole family.

\emph{Step 5: the normalizations.} The maximum of $\int_Xw\,\mathrm{d}\nu$ over $\mathcal P_f$ is attained, because under \eqref{eq:selection-face-representation} it is the maximum of the weakly continuous functional $\eta\mapsto A^{-1}\int_Xaw\eta\,\mathrm{d}V$ on $\mathcal H_a$; note that $aw\in L^\ell$. If the maximum were negative, a small positive constant translate of $w$ would still satisfy \eqref{eq:selection-class}, contrary to greatestness. Thus it is zero.

Finally, $k_\lambda\to\sup_Xw$. Indeed, from any subsequence choose a further one with $k_\lambda\to k$. Then $v_\lambda=w_\lambda-k_\lambda$ converges in $L^1$ to $w-k$, and \cref{thm:normalized-compactness} gives $\sup_X(w-k)=0$. Thus $k=\sup_Xw$ for every cluster point of $(k_\lambda)$. Subtracting the constants proves the convergence of $u_\lambda-\sup_Xu_\lambda$. If $\theta$ is Hermitian, we may take $\rho=0$ and $\Omega=X$, so that local uniform convergence is uniform convergence.
\end{proof}

Together with \cref{prop:big-exponential-clusters}, the theorem also gives convergence of the normalized potentials in the topology of uniform convergence of differences on $\Omega$ and in every local $C^\beta$ topology stated there: compactness in those topologies and uniqueness of the $L^1$ limit force every subsequence to have the same limit. The fact that the limit is selected does not imply that the critical equation has a unique normalized solution.


\section{Differentiability and dependence on the singularity type}

Return to the semipositive big background and model reference of \cref{sec:normalization-scalar}. The first question is whether $\mathcal A_\varphi(F)$ is a singleton for every continuous $F$ under \eqref{eq:als-model-condition}. Equivalently, is $B_\varphi$ G\^ateaux differentiable everywhere on $C(X)$? We have proved this for smooth data in a Hermitian background, for continuous data in the closed K\"ahler case, and on a dense $G_\delta$ for every fixed reference above. These assertions do not settle uniqueness at every continuous density in a general Hermitian background or for a singular reference.

A second question concerns a weak linearized operator at a solution $u$ whose adjoint annihilates the supporting probability. In a Hermitian background and with zero reference, \cref{thm:gaud-current-all-adjoints} will represent every supporting probability by a determinant-weighted positive $\ddc$-closed current. Identifying such a current with a scalar multiple of the weak power $\theta_u^{n-1}$ is a further issue: the product and its compatibility with approximation must first be justified.

A third question is how the sets $\mathcal A_\varphi(F)$ vary when the prescribed singularity type changes. By \cref{cor:prescribed-dual-entropy}, bounded logarithmic densities and a common smaller-background parameter give uniform entropy and compactness even for varying model references. When the smaller-background parameters tend to zero, \cref{cor:small-reference-adjoint-limits} identifies every limit with a zero-reference supporting probability. \Cref{prop:normalization-distance-stability} gives stability under uniform convergence of the scalar functionals, even for bounded measurable logarithmic densities. Relating this condition directly to the ordinary singularities or their multiplier ideals remains a separate question.

These questions concern the supporting probabilities obtained from the scalar equation. Their role differs from that of an energy on the space of potentials: the obstruction in \cref{chap:energy} still applies to a primitive of the Monge--Amp\`ere one-form.
